\chapter{Inclusion--exclusion and generating functions}\label{ch:piegf}
\index{inclusion--exclusion}\index{generating function}

Chapter~8 counted objects by naming them. This chapter handles the two
situations where naming fails: when the condition is a \emph{union}
of overlapping bad events (count the complement by
inclusion--exclusion), and when the constraint is a \emph{sum}
(track it as an exponent in a generating function, where multiplication
does the bookkeeping). After this chapter you can derive derangements
and surjection counts from one formula, extract $\varphi$ as a PIE
special case, and turn ``how many ways to make $n$'' into a coefficient
extraction.

\section{Two sets, three sets, then the formula}\label{sec:pie}
\index{principle of inclusion--exclusion}

$|A \cup B| = |A| + |B| - |A\cap B|$: the overlap was counted twice, so
subtract it once. For three sets, pairwise subtractions over-remove the
triple overlap, which is added back (Figure~\ref{fig:pievenn}).

\begin{theorem}[Inclusion--exclusion]
For finite sets $A_1, \dots, A_n$,
\[
\Bigl|\bigcup_{i=1}^{n} A_i\Bigr|
= \sum_i |A_i| - \sum_{i<j} |A_i \cap A_j|
  + \sum_{i<j<k} |A_i \cap A_j \cap A_k| - \cdots
  + (-1)^{n+1} |A_1 \cap \cdots \cap A_n|.
\]
\end{theorem}

\begin{example}\label{ex:pieproof}
Prove the general PIE formula.
\par\smallskip
\textbf{Strategy.} Show every element's total contribution is exactly
$+1$: an element in $t$ of the sets contributes
$\binom{t}{1} - \binom{t}{2} + \binom{t}{3} - \cdots = 1 - (1-1)^t = 1$ ---
algebra as bookkeeping. (The indicator-function version: expand
$1 - \prod_i (1 - \mathbb{1}_{A_i})$.)
\par\smallskip
\begin{solution}
Fix an element $x$ lying in exactly $t \ge 1$ of the sets. In the
$j$-th alternating sum it is counted $\binom{t}{j}$ times, so its total
weight is $\sum_{j \ge 1} (-1)^{j+1}\binom{t}{j}
= 1 - \sum_{j=0}^{t}(-1)^j\binom{t}{j} = 1 - (1-1)^t = 1$. Every
element contributes exactly once, hence the formula.
\end{solution}
\end{example}

\begin{figure}[ht]\centering
% Figure: 3-set Venn with +/- signs showing PIE's alternating correction.
% Source for Figure~\ref{fig:pievenn}. TikZ only; regenerable by compiling main.tex.
\begin{tikzpicture}[
  lab/.style={font=\footnotesize\sffamily,text=slate},
  reg/.style={font=\footnotesize\bfseries},
]
\begin{scope}
  \draw[draw=rulegray,fill=accent!8] (0,0) circle (1.15);
  \draw[draw=rulegray,fill=accent!8] (1.3,0) circle (1.15);
  \draw[draw=rulegray,fill=accent!8] (0.65,1.05) circle (1.15);
\end{scope}
\node[lab] at (-1.0,1.35) {$A$};
\node[lab] at (2.3,1.35) {$B$};
\node[lab] at (0.65,2.35) {$C$};
% region signs: single = +, pair = -, triple = +
\node[reg] at (-0.55,0.05) {$+$};
\node[reg] at (1.85,0.05) {$+$};
\node[reg] at (0.65,1.85) {$+$};
\node[reg] at (0.65,-0.55) {$-$};
\node[reg] at (-0.25,0.85) {$-$};
\node[reg] at (1.55,0.85) {$-$};
\node[reg] at (0.65,0.55) {$+$};
\node[lab] at (0.65,-1.75) {each element is counted $+1$ exactly once: singles $+$, pairs $-$, triple $+$};
\end{tikzpicture}

\caption{PIE's alternating correction on three sets. Notice the sign of
each region equals the parity of how many sets cover it --- the formula
is not subtracting ``mistakes'' but assigning each element weight $+1$.}
\label{fig:pievenn}
\end{figure}

\section{Derangements}\label{sec:derange}
\index{derangement}

\begin{example}\label{ex:derange}
Count derangements of $4$: permutations with no fixed point.
\par\smallskip
\textbf{Strategy.} Count the complement: let $A_i$ = permutations
fixing $i$. Then $|A_{i_1} \cap \cdots \cap A_{i_k}| = (n-k)!$ for every
$k$-fold intersection --- the bad events are easy to intersect, which is
the whole reason PIE works here.
\par\smallskip
\begin{solution}
Figure~\ref{fig:piebar} shows the bar stack:
$D_4 = 4! - \binom{4}{1}3! + \binom{4}{2}2! - \binom{4}{3}1!
+ \binom{4}{4}0! = 24 - 24 + 12 - 4 + 1 = 9$.
Generally $D_n = n!\sum_{k=0}^{n} \frac{(-1)^k}{k!} \approx n!/e$ ---
about a third of all permutations are derangements.
\end{solution}
\end{example}

\begin{figure}[ht]\centering
% Figure: PIE bar-stack for D_4 -- sizes of the correcting terms.
% Source for Figure~\ref{fig:piebar}. TikZ only; regenerable by compiling main.tex.
\begin{tikzpicture}[
  bl/.style={font=\footnotesize\sffamily,text=slate},
  val/.style={font=\footnotesize\bfseries},
]
% bars: 24, -6x4=-24? no: C(4,1)*3!=24, C(4,2)*2!=12, C(4,3)*1!=4, C(4,4)*0!=1
\draw[fill=rulegray!40] (0,0) rectangle (2.4,0.5);
\node[val] at (1.2,0.25) {$4! = 24$};
\node[bl] at (-1.1,0.25) {all};
\draw[fill=accent!35] (0,-0.9) rectangle (2.4,-0.4);
\node[val] at (1.2,-0.65) {$\binom{4}{1}\cdot 3! = 24$};
\node[bl] at (-1.1,-0.65) {$-$ singles};
\draw[fill=rulegray!40] (0,-1.8) rectangle (1.2,-1.3);
\node[val] at (0.6,-1.55) {$\binom{4}{2}\cdot 2! = 12$};
\node[bl] at (-1.1,-1.55) {$+$ pairs};
\draw[fill=accent!35] (0,-2.7) rectangle (0.4,-2.2);
\node[val] at (1.05,-2.45) {$\binom{4}{3}\cdot 1! = 4$};
\node[bl] at (-1.1,-2.45) {$-$ triples};
\draw[fill=rulegray!40] (0,-3.6) rectangle (0.1,-3.1);
\node[val] at (0.85,-3.35) {$\binom{4}{4}\cdot 0! = 1$};
\node[bl] at (-1.1,-3.35) {$+$ all};
\node[bl] at (3.1,-1.8) {$24 - 24 + 12 - 4 + 1 = 9 = D_4$};
\node[bl,anchor=west] at (0,-4.4) {fixed-point sets shrink $k!$-wise; the signs alternate};
\end{tikzpicture}

\caption{PIE for $D_4$ as a bar stack. Notice the $k$-fold intersections
shrink factorially while their counts $\binom{4}{k}$ grow --- the two
effects nearly cancel at the top, which is why derangements stay a
constant fraction of $n!$.}
\label{fig:piebar}
\end{figure}

\section{Surjections}\label{sec:surj}
\index{surjection!counting}
Same machine, different bad events: onto maps $[n] \to [k]$ are counted
by excluding the maps that miss a chosen point. $A_i$ = maps missing
target $i$; then $|A_{i_1}\cap\cdots\cap A_{i_j}| = (k-j)^n$, giving
\[
\#\{\text{surjections}\} = \sum_{j=0}^{k} (-1)^j \binom{k}{j} (k-j)^n
\;=\; \sum_{j=0}^{k} (-1)^{k-j} \binom{k}{j} j^n .
\]
(The second form reindexes $j \mapsto k-j$; the sign it carries is
$(-1)^{k-j}$, not $(-1)^j$.)

\section{Euler's totient by PIE}\label{sec:phipie}
\index{Euler's totient!by inclusion--exclusion}

\begin{example}\label{ex:phipie}
Compute $\varphi(360)$ by PIE.
\par\smallskip
\textbf{Strategy.} Coprimality is ``no prime factor shared'' --- the bad
events are the primes $2, 3, 5$ of $360 = 2^3 \cdot 3^2 \cdot 5$:
exclude multiples of each, add back pairwise multiples.
\par\smallskip
\begin{solution}
$\varphi(360) = 360 - \frac{360}{2} - \frac{360}{3} - \frac{360}{5}
+ \frac{360}{6} + \frac{360}{10} + \frac{360}{15} - \frac{360}{30}
= 360 - 180 - 120 - 72 + 60 + 36 + 24 - 12 = 96$.
This matches Chapter~7's product formula $360 \cdot \frac{1}{2} \cdot
\frac{2}{3} \cdot \frac{4}{5} = 96$ --- PIE \emph{is} the proof of that
formula: the alternating expansion factors as
$n \prod_{p \mid n}\bigl(1 - \frac{1}{p}\bigr)$.
\end{solution}
\end{example}

\section{Ordinary generating functions}\label{sec:ogf}
\index{generating function!ordinary}\index{coefficient extraction}

\begin{definition}
The ordinary generating function of a sequence $(a_n)$ is the formal
series $A(x) = \sum_{n \ge 0} a_n x^n$; $[x^n]A(x)$ denotes the
coefficient $a_n$.
\end{definition}

``Formal'' is the licence: $x$ is never evaluated, so convergence is
irrelevant --- the series is a coat-hanger for the coefficients, and
$[x^n]$ is extraction, not evaluation.

\section{The catalog and the convolution rule}\label{sec:catalog}
\index{convolution!of generating functions}
Workhorses: $\frac{1}{1-x} = \sum x^n$; $\frac{1}{1-ax} = \sum a^n x^n$;
$\frac{1}{(1-x)^2} = \sum (n+1)x^n$; $\frac{1}{1-x^k}$ counts multiples
of $k$; $(1+x)^n = \sum \binom{n}{k}x^k$ (the binomial theorem is a GF).
The key move: \textbf{multiplication convolves} --- if $C = A \cdot B$
then $[x^n]C = \sum_k a_k b_{n-k}$, which is exactly the sum rule
applied to splitting a total between two independent parts.

\section{GFs for counting}\label{sec:gfcount}

\begin{example}\label{ex:coin}
Count ways to make $30$ with coins of $5, 10, 25$ (order irrelevant,
unlimited supply).
\par\smallskip
\textbf{Strategy.} Model each coin type as the geometric factor
$1 + x^c + x^{2c} + \cdots = \frac{1}{1-x^c}$ --- choosing how many of
that coin. The product's coefficient $[x^{30}]$ counts triples
$(a,b,c)$ with $5a + 10b + 25c = 30$: the convolution does the
bookkeeping.
\par\smallskip
\begin{solution}
$\frac{1}{(1-x^5)(1-x^{10})(1-x^{25})}
= 1 + x^5 + 2x^{10} + 2x^{15} + 3x^{20} + 4x^{25} + 5x^{30} + \cdots$
(Figures~\ref{fig:gfflow}, \ref{fig:cointable}), so $[x^{30}] = 5$.
Listing~\ref{lst:gf} computes this by truncated-series multiplication
and brute-force verifies it.
\end{solution}
\end{example}

\begin{figure}[ht]\centering
% Figure: GF coefficient-extraction flowchart: factor -> expand -> read [x^n].
% Source for Figure~\ref{fig:gfflow}. TikZ only; regenerable by compiling main.tex.
\begin{tikzpicture}[
  bx/.style={draw=rulegray,rounded corners=3pt,inner sep=4pt,align=center,font=\footnotesize},
  ar/.style={->,thick,draw=ink},
  lab/.style={font=\footnotesize\sffamily,text=slate},
]
\node[bx,text width=2.6cm] (f1) at (0,0) {each choice\\$= $ a geometric\\factor};
\node[bx,text width=2.6cm] (f2) at (3.4,0) {multiply the\\series formally};
\node[bx,text width=2.6cm] (f3) at (6.8,0) {read off the\\coefficient $[x^n]$};
\draw[ar] (f1) -- (f2);
\draw[ar] (f2) -- (f3);
\node[lab] at (0,-1.2) {$\dfrac{1}{1-x^{5}}\cdot\dfrac{1}{1-x^{10}}\cdot\dfrac{1}{1-x^{25}}$};
\node[lab] at (3.4,-1.2) {$1+x^{5}+2x^{10}+2x^{15}+3x^{20}+4x^{25}+5x^{30}+\cdots$};
\node[lab] at (6.8,-1.2) {$[x^{30}] = 5$};
\node[lab] at (3.4,-2.0) {the coefficient \emph{is} the count --- $x$ never gets a value};
\end{tikzpicture}

\caption{Coefficient extraction as a pipeline. Notice $x$ is a formal
placeholder --- the coefficient \emph{is} the count, and no value is
ever substituted.}
\label{fig:gfflow}
\end{figure}

\begin{figure}[ht]\centering
% Figure: coin-GF partial-product table (coefficients through x^30).
% Source for Figure~\ref{fig:cointable}. Tabular; regenerable by compiling main.tex.
\begin{tabular}{c|rrrrrrr}
\toprule
factor & $[x^0]$ & $[x^5]$ & $[x^{10}]$ & $[x^{15}]$ & $[x^{20}]$ & $[x^{25}]$ & $[x^{30}]$\\
\midrule
$1/(1-x^5)$            & 1 & 1 & 1 & 1 & 1 & 1 & 1\\
$\times\,1/(1-x^{10})$ & 1 & 1 & 2 & 2 & 3 & 3 & 4\\
$\times\,1/(1-x^{25})$ & 1 & 1 & 2 & 2 & 3 & 4 & \textbf{5}\\
\bottomrule
\end{tabular}

\caption{Partial products for the coin GF through $x^{30}$. Notice each
new factor adds the previous row shifted by the coin's step --- the
table \emph{is} the convolution, run by hand.}
\label{fig:cointable}
\end{figure}

\begin{checkpoint}[Predict before you run]
Before looking at Listing~\ref{lst:gf}'s output: predict whether
$[x^{30}]$ for coins $\{5,10,25\}$ is odd or even, and name one of the
five solutions.
\end{checkpoint}

\begin{lstlisting}[language=Python,caption={\texttt{gf\_coin.py}: GF
coefficients by truncated-series multiplication + brute-force check
(script in \texttt{code/gf\_coin.py}).},label={lst:gf}]
for c in coins:
    ng = g[:]            # convolve with the coin's series
    for k in range(c, n + 1, c):
        for j in range(n - k + 1):
            ng[j + k] += g[j]
\end{lstlisting}
Real output:
\begin{lstlisting}[numbers=none]
GF product for coins {5,10,25}, terms through x^30:
  coefficients [x^0..x^30] nonzero: {0: 1, 5: 1, 10: 2, 15: 2, 20: 3, 25: 4, 30: 5}
[x^30] by GF:  5
[x^30] by brute force: 5
\end{lstlisting}

Constrained integer solutions are the same shape:
$x_1 + x_2 + x_3 = n$ with $x_i \le c_i$ is $[x^n]$ of
$\prod_i (1 + x + \cdots + x^{c_i}) = \prod_i \frac{1-x^{c_i+1}}{1-x}$ ---
the stars-and-bars count corrected by PIE, with the correction factors
already built into the numerator.

\section{GFs meet recurrences}\label{sec:gfrec}
\index{generating function!and recurrences}
A recurrence is a GF equation: Fibonacci $f_{n} = f_{n-1} + f_{n-2}$
translates to $F(x) = x + xF(x) + x^2F(x)$, hence
$F(x) = \frac{x}{1-x-x^2}$ --- coefficients recoverable by partial
fractions. That is Chapter~10's machinery; the point here is the
conversion habit: recurrence in, rational function out, coefficients
back.

\begin{exercise}[easy]\label{exr:pie2}
How many integers in $1..1000$ are divisible by 3 or 5? By neither?
\end{exercise}
\begin{exercise}[medium]\label{exr:d5}
Compute $D_5$ from the general formula; check that $D_5/5!$ is closer
to $1/e$ than $D_4/4!$ was.
\end{exercise}
\begin{exercise}[medium]\label{exr:gf}
What is $[x^{12}]$ of $\frac{1}{(1-x^2)(1-x^3)}$? Interpret it as a
coin count.
\end{exercise}
\begin{exercise}[medium]\label{exr:surj}
Count surjections $[4] \to [3]$ and verify against $3^4$ minus the
miss-a-point cases by hand.
\end{exercise}

\section{Chapter summary}
\begin{table}[ht]\centering
\begin{tabularx}{\linewidth}{@{}lYY@{}}
\toprule
\tablehead\thead{Want} & \thead{Tool} & \thead{Shape}\\
\midrule
union of bad events & PIE & alternating intersections\\
no fixed point & derangement & $n!\sum(-1)^k/k!$\\
onto maps $[n]\to[k]$ & surjection count & $\sum(-1)^j\binom{k}{j}j^n$\\
coprimes in $[n]$ & $\varphi$ via PIE & $n\prod_{p|n}(1-1/p)$\\
constrained sums & generating function & $[x^n]\prod(\cdot)$\\
recurrences & GF algebra & solve for $F$, extract $[x^n]$\\
\bottomrule
\end{tabularx}
\caption{One page: bad events $\to$ PIE; sum constraints $\to$
coefficient extraction; recurrences $\to$ GF algebra.}
\label{tab:piegf}
\end{table}

\section{Common mistakes}
\index{common mistakes}
\begin{itemize}
\item \textbf{Stopping PIE at single sets.} $|A\cup B\cup C| \ne
|A|+|B|+|C|$ even when pairs are subtracted --- the triple overlap must
come back. Sign discipline: subtract intersections of \emph{even}
order, add \emph{odd} order, and check a single element's weight is
$+1$.
\item \textbf{Treating $[x^n]$ as evaluation.} $[x^{30}]$ is the
coefficient of the series, not the value at $x = 30$. GFs are formal:
$x$ is a bookkeeping placeholder, never a number.
\item \textbf{Reading the wrong coefficient.} After multiplying
factors, the answer lives at the \emph{combined} exponent ---
$[x^{n}]$ of the product, where $n$ is the total being made, not at any
factor's own exponent.
\end{itemize}

\section{Additional exercises}
\begin{exercise}[easy]
Of 200 students, 120 take maths, 90 take physics, 50 take both. How
many take neither?
\end{exercise}
\begin{exercise}[easy]
Compute $\varphi(1000)$ by PIE (prime factors $2, 5$) and check against
Chapter~7's product formula.
\end{exercise}
\begin{exercise}[easy]
Extract $[x^6]$ from $\frac{1}{1-x} \cdot \frac{1}{1-x^3}$ and name
the counting problem it answers.
\end{exercise}
\begin{exercise}[easy]
List the 9 derangements of 4 elements explicitly (or write the
generator) and verify a sample of three fix nothing.
\end{exercise}
\begin{exercise}[medium]
Count ways to make $40$ with coins $\{5,10,25\}$ using the partial-
product method; verify against \texttt{gf\_coin.py}.
\end{exercise}
\begin{exercise}[medium]
Prove the derangement recurrence $D_n = (n-1)(D_{n-1} + D_{n-2})$
directly (where does element 1 go?), then reconcile with the PIE
formula.
\end{exercise}
\begin{exercise}[medium]
Show $\varphi(n) = n\prod_{p\mid n}(1-1/p)$ by PIE over the distinct
prime factors --- why do only the \emph{distinct} primes enter?
\end{exercise}
\begin{exercise}[medium]
Count solutions of $x_1+x_2+x_3 = 12$ with each $x_i \le 5$, once by
PIE-on-stars-and-bars and once as $[x^{12}](1+x+\cdots+x^5)^3$.
\end{exercise}
\begin{exercise}[hard]
Prove the surjection formula $\sum_j(-1)^j\binom{k}{j}j^n$ gives $0$
when $n < k$ --- explain combinatorially why it \emph{must}.
\end{exercise}
\begin{exercise}[hard]
Derive the GF for partitions into distinct parts and into odd parts,
and prove the two counts are equal by manipulating the products.
\end{exercise}
\begin{exercise}[hard]
Extend \texttt{gf\_coin.py} to print the coefficient table for any coin
list and any $n$, then use it to find the largest amount that cannot be
made with $\{6, 10, 15\}$ (answer: $29$; prove it is largest by
induction on residues mod 5 --- and first check why the question only
makes sense when $\gcd = 1$).
\end{exercise}

\chapternotes
PIE's indicator-function proof and the derangement treatment follow
Hammack's counting chapter~\cite{hammack} and 6.042J's counting
sections~\cite{mcs}; the GF-as-bookkeeping viewpoint is CS70's
combinatorial-proofs note~\cite{cs70} and Chen's summation handout's GF
section~\cite{chen-summ}; Engel supplies the hard-end drill~\cite{engel}.
Contest items powered by this chapter --- PIE counts (AIME~2024~II~P1,
AIME~2010~I~P7), derangements and surjections (AIME~2022~I~P12--13,
AIME~2015~I~P12, AIME~2018~I~P12), GF extractions (AIME~2023~II~P7,
AIME~2023~II~P11, AMC~12~2023~A~P24, AIME~2025~II~P8) --- are on the
AoPS archive~\cite{aops-amc}.
